\documentclass[11pt]{article}
\usepackage[a4paper,margin=24mm]{geometry}
\usepackage{amsmath,amssymb,amsthm}
\usepackage[T1]{fontenc}
\usepackage{lmodern}
\usepackage[hidelinks]{hyperref}
\newtheorem{lemma}{Lemma}
\newtheorem{theorem}{Theorem}
\newcommand{\AF}{\operatorname{AF}}
\newcommand{\FF}{\operatorname{FF}}
\setlength{\parindent}{0pt}
\setlength{\parskip}{5pt}
\begin{document}
\begin{center}
{\LARGE A disproof of conjecture 00000004289\par}
\vspace{5pt}
{\large Every almost-free group in the stated sense is flat-free\par}
\vspace{10pt}
C0ldSmi1e \quad\textbullet\quad October 5, 2026
\end{center}

\section{Definitions and the assertion being refuted}
Let $G$ be an abelian group, regarded as a module over $\mathbb Z$, and let
$\kappa$ be any cardinal. The bilingual source defines the properties
\begin{align*}
 \AF_\kappa(G)&\iff G\text{ is not free and every subgroup }H\leq G
                 \text{ with }|H|<\kappa\text{ is free},\\
 \FF_\kappa(G)&\iff \text{every subgroup }H\leq G
                 \text{ with }|H|<\kappa\text{ is flat}.
\end{align*}
Free means free as a $\mathbb Z$-module, with an arbitrary basis and
finite-support coordinates. Flat means that tensoring over $\mathbb Z$
preserves exact sequences, equivalently that it preserves injective module
maps; see \cite{stacks}. All subgroups are included, and the cardinal bound
is strict in both definitions. The non-free requirement belongs to
$\AF_\kappa$, as in the source.

For clarity, keep the threshold and the group size separate:
\[
 E(\kappa,\lambda)\iff
 \exists G\;\bigl(|G|=\lambda\ \text{and}\
 \AF_\kappa(G)\ \text{and}\ \neg\FF_\kappa(G)\bigr).
\]
The usual reading of ``at cardinal $\lambda$'' uses threshold
$\kappa=\lambda$. On this reading the conjecture asserts that
$\aleph_\omega$ is the least member of
$S=\{\lambda:E(\lambda,\lambda)\}$.
In particular, it asserts $E(\aleph_\omega,\aleph_\omega)$.
We disprove this necessary existence clause and the minimum claim.
The argument also covers a fixed threshold or a threshold unrelated to the
group's size.

\section{Complete proof}
\begin{lemma}
Every free abelian group is a flat $\mathbb Z$-module.
\end{lemma}
\begin{proof}
If $F$ has basis $I$, then $F\cong\bigoplus_{i\in I}\mathbb Z$.
For every abelian group $A$ there is a natural identification
$A\otimes_{\mathbb Z}F\cong\bigoplus_{i\in I}A$.
For an injective homomorphism $f:A\to B$, the induced tensor map becomes
the coordinatewise map $\bigoplus_I f$. If its value on a finite-support
tuple is zero, injectivity of $f$ makes each coordinate zero. Thus the
tensor map is injective. The injectivity criterion for flatness proves
the lemma, including an empty or infinite basis.
\end{proof}

\begin{theorem}
For every cardinal $\kappa$ and every abelian group $G$,
$\AF_\kappa(G)$ implies $\FF_\kappa(G)$.
Consequently $E(\kappa,\lambda)$ is false for all cardinals $\kappa,\lambda$.
\end{theorem}
\begin{proof}
Assume $\AF_\kappa(G)$. Let $H\leq G$ satisfy $|H|<\kappa$.
The definition makes $H$ free, so the lemma makes $H$ flat.
This proves $\FF_\kappa(G)$ and contradicts the last condition in
any proposed witness to $E(\kappa,\lambda)$.
\end{proof}

\newpage
\section{Consequences and logical scope}
There is no separating group at $\aleph_\omega$, or at any other
cardinality. The set $S$ is empty, so it has no least member.
A least element must belong to the set; a mere lower bound or an infimum
of an empty set would not express the source's asserted minimum.

The separate $\aleph_1$ implication in the source is true, by the theorem
with $\kappa=\aleph_1$, even without requiring $|G|=\aleph_1$.
No regularity, cofinality, infinitude, or positivity assumption on the
threshold is used. If the subgroup range is empty, the same implication
still holds. The proof does not require the non-free conjunct, but the
formal definition retains it exactly.

The source also asserts ZFC provability of the separation, a Shelah black-box
construction, and independence over ZFC at every smaller cardinal.
Refuting the necessary existence and minimum clauses disproves this compound
assertion. This submission does not formalize a coded ZFC proof system,
independence, or a black-box construction; a Lean theorem alone is not such
a formalization. The unspecified range of ``smaller cardinal'' does not
affect the universal nonexistence result.

\section{Formal correspondence and verification}
The project uses Lean 4.19.0 and Mathlib v4.19.0. The carrier has type
\texttt{G : Type u}. Its \texttt{AddCommGroup G} structure provides the
canonical integer action. Freeness is Mathlib's basis-existence property,
\texttt{Module.Free}. Flatness is its tensor-flatness class,
\texttt{Module.Flat}, with proved injectivity and exactness interpretations
over rings. The instance \texttt{Module.Flat.of\_free} supplies the lemma.
Subgroups are actual elements of \texttt{AddSubgroup G}, with inherited
group operations and cardinality given by \texttt{Cardinal.mk}.

\texttt{AlmostFree}, \texttt{FlatFree}, \texttt{Separation}, and
\texttt{SeparationExists} are exactly the predicates stated on page~1.
\texttt{SeparationCardinals} is the set $S$.
The theorem \texttt{no\_separation\_exists} quantifies independently over
$\kappa$ and $\lambda$ in every universe. The final theorem
\texttt{conjecture4289\_disproof} states
$\neg\operatorname{IsLeast}(S,\aleph_\omega)$.
The cardinal $\aleph_\omega$ is \texttt{Cardinal.aleph Ordinal.omega0};
$\aleph_1$ is \texttt{Cardinal.aleph 1}.
Separate theorems exclude existence at $\aleph_\omega$, show $S$ empty,
and specialize the implication at $\aleph_1$.

The proof is symbolic and needs no external mathematical computation.
From the pinned project in \texttt{lean/}, run
\begin{verbatim}
lake build
lake env lean -DwarningAsError=true Conjecture4289.lean
lake env lean -DwarningAsError=true Check.lean
\end{verbatim}
The audit prints every authored definition, theorem type, and axiom dependency.
Verification records cover all compiled proof-module constants, including
generated ones, dependency identities, and the matching PDF. Local checks and
independent internal review do not establish maintainer acceptance.

\begin{thebibliography}{9}
\bibitem{stacks} The Stacks Project, Section 10.39,
\emph{Flat modules and flat ring maps}, Definition 10.39.1 and Lemma 10.39.5.
\href{https://stacks.math.columbia.edu/tag/00H9}{stacks.math.columbia.edu/tag/00H9}.
\bibitem{mathlib} Mathlib v4.19.0,
\texttt{Mathlib/RingTheory/Flat/Basic.lean},
including \texttt{Module.Flat.of\_free}.
The exact dependency revision is recorded in the pinned project manifest.
\end{thebibliography}
\end{document}
